\chapter{Running a simulation}\label{runningMoDELib}
\section{Build the tools}
The repository uses C++20 and CMake 3.20 or later. Eigen and FFTW are required by the current build. OpenMP is used by the simulation tools. Qt/VTK are needed for the optional viewer, and pybind11 for optional Python bindings. The root \code{CMakeLists.txt} is the authoritative list of build options for this snapshot.

From the repository root, a command-line configuration is:
\par\noindent\begin{minipage}{\linewidth}
\begin{lstlisting}[language=bash]
cmake -S . -B build -DUSE_PYBIND11=OFF -DBUILD_DDQT=OFF
cmake --build build -j
\end{lstlisting}
\end{minipage}\par
This builds the simulation tools without requiring the viewer or Python bindings. Tutorial input generators still need Python and NumPy. Plotting scripts generally need Matplotlib. CD species counts are compile-time options; match them to the chosen CD material when that physics is enabled.

\section{The simulation folder}
A simulation has an \code{inputFiles/} directory, a \code{config/} directory, and an \code{F/} directory for history output. The generator writes the initial microstructure; the driver evolves it. Copy templates into the simulation before editing them.
\begin{center}\small
\begin{tabular}{p{0.35\textwidth}p{0.57\textwidth}}\toprule
File & Responsibility\\\midrule
\code{polycrystal.txt} & Mesh, transformations, temperature, material, grain orientations and periodic faces.\\
\code{DefectiveCrystal.txt} & Physics selection, shared FEM switch, time stepping, image counts and output schedule.\\
Material file & Units, elasticity, crystal structure and constitutive parameters.\\
\code{DD.txt} & DDD solver, remeshing and dislocation-specific controls.\\
\code{ElasticDeformation.txt} & Uniform/FEM elasticity, loading and rigid-motion control.\\
\code{initialMicrostructure.txt} & List of microstructure specifications consumed by the generator.\\
\code{ClusterDynamics.txt} & Current CD-specific controls, when enabled.\\\bottomrule
\end{tabular}\end{center}
Inputs generally use \code{name=value;} and \code{\#} comments. Arrays are whitespace-separated; matrices have rows and one terminating semicolon. Use the complete template for each enabled physics rather than assembling a deck from isolated examples in this manual.

\section{A first DDD run}
The \code{dipoles} tutorial demonstrates the complete path. Run in a fresh copy if that tutorial folder contains results you wish to keep. From the repository root:
\par\noindent\begin{minipage}{\linewidth}
\begin{lstlisting}[language=bash]
cd tutorials/dipoles
python3 generateInputFiles.py
../../build/tools/MicrostructureGenerator/microstructureGenerator .
../../build/tools/DDomp/DDomp .
\end{lstlisting}
\end{minipage}\par
The supplied generator selects DDD and elastic deformation, copies the material and mesh, generates \code{polycrystal.txt}, and writes the microstructure specifications. Its default is a long run. For an initial check, edit the generated \code{DefectiveCrystal.txt} to set \code{Nsteps=2;} and \code{outputFrequency=1;} before invoking \code{DDomp}. The commands here document the workflow; no simulation was run as part of this manual revision.

\section{Outputs and restart}
\code{config/evl_*.txt} (or the binary equivalent) contains configurations. Auxiliary output carries additional fields selected by the physics. \code{F/F_0.txt} stores scalar history, and \code{F/F_labels.txt} identifies its columns. Read labels instead of assuming a fixed column layout across physics combinations.

\code{outputFrequency} controls step-based output and must be positive; \code{outputBinary} selects configuration encoding. \code{startAtTimeStep=-1} requests the latest available history step. For a fresh run, ensure that old configuration and history files are not accidentally used as a restart. For continuation, retain the configuration and matching history and set \code{Nsteps} beyond the restored step index.

Check the selected physics, unit scales, periodic basis and initial configuration before interpreting results. Compare time-step, spatial and image-sum resolution using the physical observable of interest. A run that completes successfully can still be under-resolved.
\sourcefile{tutorials/dipoles/generateInputFiles.py}
\sourcefile{src/DefectiveCrystal/DefectiveCrystalParameters.cpp}
\sourcefile{src/DefectiveCrystal/DefectiveCrystalTraitsIO.cpp}

\section{Periodic mesh input}
Set \code{periodicFaceIDs} in \code{polycrystal.txt}: an empty set disables pairing; \code{periodicFaceIDs=-1;} enables all valid discovered directions. Explicit IDs select external planar MoDELib faces, not Gmsh physical surface tags. Inspect the mesh face numbering and the discovered translation basis before selecting a subset. Geometric matching is described in Section~\ref{sec:periodicity}; FEM constraints and DDD image controls are described in Sections~\ref{sec:periodicFEM} and~\ref{sec:dddPeriodicity}, respectively.
